
% !TEX program = pdflatex
% ============================================================
% Learning notes on LHmuC Higgs and BSM benchmark slides
%
% Compile locally in the folder containing:
%   LHmuC_vector.pdf
%   LHmuC_fermion.pdf
%   LHmuC_loop.pdf
%   exclusion_limits.jpg
%   mu8.jpg
%
% Suggested:
%   pdflatex LHmuC_Higgs_BSM_learning_notes.tex
%   pdflatex LHmuC_Higgs_BSM_learning_notes.tex
% ============================================================

\documentclass[11pt,a4paper]{article}

\usepackage[utf8]{inputenc}
\usepackage[T1]{fontenc}
\usepackage[a4paper,margin=2.2cm]{geometry}
\usepackage{lmodern}
\usepackage{amsmath,amssymb,bm}
\usepackage{graphicx}
\usepackage{xcolor}
\usepackage{booktabs}
\usepackage{array}
\usepackage{longtable}
\usepackage{enumitem}
\usepackage{caption}
\usepackage{subcaption}
\usepackage{float}
% ---------- Commands ----------
\newcommand{\LHmuC}{LH\ensuremath{\mu}C}
\newcommand{\muLHC}{\ensuremath{\mu}LHC}
\newcommand{\LHeC}{LHeC}
\newcommand{\HL}{HL-LHC}
\newcommand{\FCCee}{FCC-ee}
\newcommand{\mup}{\ensuremath{\mu p}}
\newcommand{\mupplus}{\ensuremath{\mu^+p}}
\newcommand{\mupminus}{\ensuremath{\mu^-p}}
\newcommand{\gev}{\ensuremath{~\mathrm{GeV}}}
\newcommand{\tev}{\ensuremath{~\mathrm{TeV}}}
\newcommand{\invfb}{\ensuremath{~\mathrm{fb}^{-1}}}
\newcommand{\invab}{\ensuremath{~\mathrm{ab}^{-1}}}
\newcommand{\pb}{\ensuremath{~\mathrm{pb}}}
\newcommand{\fb}{\ensuremath{~\mathrm{fb}}}
\newcommand{\BR}{\ensuremath{\mathrm{BR}}}
\newcommand{\ghbb}{\ensuremath{g_{hbb}}}
\newcommand{\ghWW}{\ensuremath{g_{hWW}}}
\newcommand{\Gamh}{\ensuremath{\Gamma_h}}
\newcommand{\kW}{\ensuremath{\kappa_W}}
\newcommand{\kZ}{\ensuremath{\kappa_Z}}
\newcommand{\kV}{\ensuremath{\kappa_V}}
\newcommand{\muoct}{\ensuremath{\mu_8}}

% ---------- Safe figure inclusion ----------
\newcommand{\includefig}[2][]{%
  \IfFileExists{#2}{\includegraphics[#1]{#2}}{%
  \fbox{\begin{minipage}[c][3.5cm][c]{0.85\linewidth}
  \centering Missing figure file: \texttt{#2}
  \end{minipage}}}%
}

% ---------- Lightweight note boxes ----------
\definecolor{aghgreen}{RGB}{0,105,65}
\definecolor{aghred}{RGB}{170,20,50}
\definecolor{darkblue}{RGB}{30,60,110}

\newenvironment{keypoint}[1]{%
\par\medskip\noindent{\large\bfseries\color{aghgreen}#1}\par\noindent
\begin{minipage}{0.97\linewidth}\itshape
}{\end{minipage}\par\medskip}

\newenvironment{warningnote}[1]{%
\par\medskip\noindent{\large\bfseries\color{aghred}#1}\par\noindent
\begin{minipage}{0.97\linewidth}
}{\end{minipage}\par\medskip}

\newenvironment{learningnote}[1]{%
\par\medskip\noindent{\large\bfseries\color{darkblue}#1}\par\noindent
\begin{minipage}{0.97\linewidth}
}{\end{minipage}\par\medskip}

\title{\vspace{-1.0cm}
\textbf{Higgs Precision and BSM Benchmarks at \LHmuC}\\[1mm]
\large Charged-current \(h\to b\bar b\), \(\kappa\)-framework fits, RPV squarks, and color-octet muons}

%\author{Prepared as reading notes for the LH\(\mu\)C / HELP4CERN presentation}

\date{\today}

\begin{document}
\maketitle

\begin{abstract}
These notes collect the physics discussion behind four presentation topics:
(i) the precision Higgs case study \(p\mu^+\to j\bar\nu_\mu h\), \(h\to b\bar b\);
(ii) the broader Higgs-coupling reach in a \(\kappa\)-framework global-fit picture;
(iii) why \LHmuC{} is also a distinctive BSM machine; and
(iv) representative BSM benchmarks, namely R-parity-violating squarks and color-octet muons.
The emphasis is pedagogical: what the slides are claiming, how the numbers are obtained,
what the plots mean, and what caveats should be stated during a conference talk.
\end{abstract}



\section{Executive summary}



\begin{keypoint}{The short physics story}
\vspace{0.25cm}
\LHmuC{} is not simply an energy upgrade of the LHeC.  It opens a multi-TeV
lepton--hadron regime where weak-boson-fusion Higgs production, high-\(Q^2\)
electroweak measurements, photon-induced reactions, top production, and direct
\(\mu q/\mu g\) BSM resonances become quantitatively competitive.
\end{keypoint}


The four target slides should be understood as one connected narrative:

\begin{enumerate}[leftmargin=*]
  \item \textbf{Concrete Higgs benchmark:} charged-current VBF production,
  \(p\mu^+\to j\bar\nu_\mu h\), followed by \(h\to b\bar b\), can be translated into
  a sensitivity on \(\ghbb\).
  \item \textbf{Global Higgs fit:} beyond a single channel, the strongest visible improvement
  is in vector modifiers, especially \(\kW\) and \(\kZ\), because Higgs production at
  \LHmuC{} is dominated by weak-boson fusion.
  \item \textbf{BSM logic:} \LHmuC{} is special because it directly probes \(\mu q\) and
  \(\mu g\) initial states at multi-TeV energy.
  \item \textbf{BSM benchmarks:} RPV squarks test the \(\mu q\) portal, while color-octet
  muons test the \(\mu g\) portal.
\end{enumerate}

\section{Benchmark setup and notation}

The current presentation uses the following working points:
\begin{table}[H]
\centering
\caption{Nominal \LHmuC{} benchmark points used in the slides.}
\begin{tabular}{@{}lccc@{}}
\toprule
Scenario & \(E_\mu\) & \(E_p\) & \(\sqrt{s_{\mu p}}\) and luminosity \\
\midrule
\LHmuC{} baseline & \(500\gev\) & \(7\tev\) &
\(\sqrt{s}\simeq3.7\tev,\quad L_{\rm int}=250\invfb\)\\
\LHmuC{} upgrade & \(1\tev\) & \(7\tev\) &
\(\sqrt{s}\simeq5.3\tev,\quad L_{\rm int}=1000\invfb\)\\
Possible higher stage & \(2\tev\) & \(7\tev\) &
\(\sqrt{s}\simeq7.5\tev,\quad L_{\rm int}=1000\invfb\)\\
\bottomrule
\end{tabular}
\end{table}

\begin{warningnote}{Important consistency issue}
The \(h\to b\bar b\) table in the slide gives
\(\Delta\sigma/\sigma=0.69\%\) for the \(1\tev\) \LHmuC{} benchmark.
Using the table cross sections and efficiencies, this value corresponds more closely to
\(L_{\rm int}=500\invfb\), not \(1000\invfb\).  If the slide keeps the global benchmark
\(L_{\rm int}=1000\invfb\), the rate uncertainty should be closer to \(0.49\%\), under
the same assumptions.  This point should be fixed before the talk.
\end{warningnote}



\section{Precision Higgs case study: \(h\to b\bar b\)}

\subsection{Why charged-current VBF dominates}


At a lepton--hadron collider, Higgs production is mainly electroweak vector-boson fusion:
\begin{align}
  \text{CC VBF:}\quad & p\mu^+\to j\,\bar\nu_\mu\,h,\\
  \text{NC VBF:}\quad & p\mu^+\to j\,\mu^+\,h.
\end{align}
The charged-current process is larger because it is driven by \(W\)-exchange.  In the current
benchmarks, the slides quote approximately
\begin{align}
  \sigma_{\rm CC}(500\gev)&\simeq438\fb, &
  \sigma_{\rm NC}(500\gev)&\simeq135\fb,\\
  \sigma_{\rm CC}(1\tev)&\simeq731\fb, &
  \sigma_{\rm NC}(1\tev)&\simeq221\fb.
\end{align}
This already gives large Higgs statistics:
\begin{align}
  N_{\rm CC}(500\gev)&\simeq 438\fb\times250\invfb \simeq1.10\times10^5,\\
  N_{\rm CC}(1\tev)&\simeq 731\fb\times1000\invfb \simeq7.31\times10^5.
\end{align}

\begin{learningnote}{Why the muon beam matters}
\vspace{0.25cm}
The physics is analogous to the LHeC, but a muon beam allows much higher lepton energy
because synchrotron radiation is much weaker for muons than electrons.  This raises
\(\sqrt{s_{\mu p}}\), increases the hard-scattering phase space, and enhances Higgs, top,
electroweak, and BSM rates.
\end{learningnote}

\subsection{Why \(h\to b\bar b\) is the natural benchmark}

The decay \(h\to b\bar b\) is the largest Standard Model Higgs decay mode and directly
probes the bottom Yukawa coupling.  It is experimentally nontrivial because the final state
contains hadronic jets and backgrounds from \(Z\to b\bar b\) and multijet production.
However, the lepton--hadron environment is cleaner than \(pp\) collisions and the topology
is more constrained.

The benchmark signal is
\begin{equation}
  p\mu^+\to j\,\bar\nu_\mu\,h,\qquad h\to b\bar b .
\end{equation}
The dominant backgrounds are
\begin{align}
  p\mu^+&\to j\,\bar\nu_\mu\,Z,\qquad Z\to b\bar b,\\
  p\mu^+&\to \bar\nu_\mu\,j\,j\,j .
\end{align}



\subsection{Analysis chain}

The benchmark analysis uses
\begin{itemize}[leftmargin=*]
  \item MG5\_aMC@NLO for event generation;
  \item Pythia6 for showering and hadronization;
  \item FastJet for jet clustering;
  \item Delphes with an LHeC-like asymmetric detector card;
  \item \(75\%\) \(b\)-tagging efficiency;
  \item exactly two reconstructed \(b\)-jets;
  \item the Higgs mass-window cut
  \begin{equation}
    |m_{bb}-m_h|<25\gev.
  \end{equation}
\end{itemize}



\subsection{Rate and coupling interpretation}

The measured final-state rate is
\begin{equation}
  \sigma(p\mu^+\to j\bar\nu_\mu b\bar b)
  =
  \sigma(p\mu^+\to j\bar\nu_\mu h)\,
  {\rm BR}(h\to b\bar b).
\end{equation}
Since charged-current VBF scales as \(g_{hWW}^2\), while
\[
  {\rm BR}(h\to b\bar b)\propto \frac{g_{hbb}^2}{\Gamma_h},
\]
the final-state rate scales as
\begin{equation}
  \sigma(p\mu^+\to j\bar\nu_\mu b\bar b)
  \propto
  \frac{g_{hWW}^2\,g_{hbb}^2}{\Gamma_h}.
  \label{eq:hbb_rate}
\end{equation}
Thus the \(hbb\) channel does not measure \(g_{hbb}\) alone.  It measures a product involving
the production coupling, the decay coupling, and the total Higgs width.

The event yields are
\begin{align}
  N_s &= L_{\rm int}\,\sigma_{\rm sig}\,\epsilon_{\rm sig},\\
  N_b &= L_{\rm int}\left[
    \sigma_{\rm bgd-Z}\epsilon_{\rm bgd-Z}
    +
    \sigma_{\rm bgd-mj}\epsilon_{\rm bgd-mj}
  \right].
\end{align}
The statistical rate uncertainty is
\begin{equation}
  \frac{\Delta\sigma}{\sigma}
  =
  \frac{\sqrt{N_s+N_b}}{N_s}.
  \label{eq:rate_unc}
\end{equation}

To extract \(\ghbb\), external information on \(g_{hWW}^2/\Gamma_h\) is used:
\begin{equation}
  \frac{\Delta g_{hbb}}{g_{hbb}}
  =
  \frac{1}{2}
  \sqrt{
  \left(\frac{\Delta\sigma}{\sigma}\right)^2+
  \left[
  \frac{
  \Delta(g_{hWW}^2/\Gamma_h)
  }{
  g_{hWW}^2/\Gamma_h
  }
  \right]^2
  }.
  \label{eq:ghbb_precision}
\end{equation}
The factor \(1/2\) appears because the rate depends on \(g_{hbb}^2\), not linearly on
\(g_{hbb}\).





\subsection{Numerical check}


For the \(500\gev\) benchmark, using
\[
  L_{\rm int}=250\invfb=250000\pb^{-1},
\]
\[
  \sigma_{\rm sig}=0.24\pb,\quad \epsilon_{\rm sig}=0.22,
\]
\[
  \sigma_{\rm bgd-Z}=0.30\pb,\quad \epsilon_{\rm bgd-Z}=0.025,
\]
\[
  \sigma_{\rm bgd-mj}=357.9\pb,\quad \epsilon_{\rm bgd-mj}=9.7\times10^{-5},
\]
one obtains
\begin{align}
  N_s &\simeq 250000\times0.24\times0.22\simeq13200,\\
  N_{b,Z} &\simeq 250000\times0.30\times0.025\simeq1875,\\
  N_{b,{\rm mj}} &\simeq 250000\times357.9\times9.7\times10^{-5}\simeq8680.
\end{align}
Then
\begin{equation}
  \frac{\Delta\sigma}{\sigma}
  \simeq
  \frac{\sqrt{13200+1875+8680}}{13200}
  \simeq1.17\%,
\end{equation}
consistent with the quoted \(1.19\%\).

For the \(1\tev\) benchmark, using
\[
  \sigma_{\rm sig}=0.39\pb,\quad \epsilon_{\rm sig}=0.21,\quad
  \sigma_{\rm bgd-Z}=0.50\pb,\quad \epsilon_{\rm bgd-Z}=0.025,
\]
\[
  \sigma_{\rm bgd-mj}=563.8\pb,\quad \epsilon_{\rm bgd-mj}=1.23\times10^{-4},
\]
one obtains
\[
  \Delta\sigma/\sigma\simeq0.49\%
\]
if \(L_{\rm int}=1000\invfb\), but
\[
  \Delta\sigma/\sigma\simeq0.69\%
\]
if \(L_{\rm int}=500\invfb\).  This is the origin of the consistency warning above.

Using the external uncertainty
\[
  \frac{\Delta(g_{hWW}^2/\Gamma_h)}{g_{hWW}^2/\Gamma_h}=0.985\%,
\]
the quoted coupling precisions are reproduced:
\begin{align}
  \frac{\Delta g_{hbb}}{g_{hbb}}
  &=
  \frac12\sqrt{(1.19\%)^2+(0.985\%)^2}
  \simeq0.77\%,\\
  \frac{\Delta g_{hbb}}{g_{hbb}}
  &=
  \frac12\sqrt{(0.69\%)^2+(0.985\%)^2}
  \simeq0.60\%.
\end{align}



\section{Broader Higgs-coupling reach: global-fit picture}

\subsection{The \(\kappa\)-framework}


In the \(\kappa\)-framework, Higgs couplings are rescaled relative to the SM:
\begin{equation}
  g_{hXX}=\kappa_X\,g_{hXX}^{\rm SM}.
\end{equation}
The relevant modifiers include:
\begin{align}
  \kappa_W,\ \kappa_Z &: \text{couplings to weak vector bosons},\\
  \kappa_b,\ \kappa_c,\ \kappa_t,\ \kappa_\tau,\ \kappa_\mu &: \text{fermion Yukawa modifiers},\\
  \kappa_g,\ \kappa_\gamma,\ \kappa_{Z\gamma} &: \text{effective loop-induced modifiers}.
\end{align}
Additional BSM branching fractions can be included:
\[
  {\rm Br}_{\rm inv},\qquad {\rm Br}_{\rm und}.
\]

A measured signal strength behaves schematically as
\begin{equation}
  \mu(i\to h\to f)
  \sim
  \frac{\kappa_i^2\kappa_f^2}{\kappa_h^2},
\end{equation}
where \(\kappa_h^2\) encodes the total-width rescaling.  This is why a global fit is needed:
individual channels measure products of production and decay couplings divided by the total width.



\subsection{The \(\kappa_V\le1\) condition}

The condition

\begin{equation}
  \kappa_V\le1,\qquad V=W,Z,
\end{equation}
is often imposed in future-collider Higgs fits when the total Higgs width is not directly and
model-independently measured.  It helps break a flat direction between the total width and
the couplings.  This is a useful benchmark assumption, but it is not fully model independent;
some BSM scenarios can violate it.



\begin{warningnote}{How to phrase this in a talk}
\vspace{0.25cm}
Say: ``The fit uses the standard \(\kappa_V\le1\) benchmark assumption to control the
Higgs-width flat direction.  This is useful for comparisons, but it is not a fully
model-independent statement.''
\end{warningnote}

\subsection{How the \LHmuC{} projection is obtained}

The global-fit result in the slide should be described as an \textbf{exploratory projection}.
The idea is:
\begin{enumerate}[leftmargin=*]
  \item start from LHeC Higgs signal-strength projections;
  \item keep the LHeC systematic scale factors fixed;
  \item rescale statistical uncertainties using \LHmuC{} benchmark cross sections and luminosities;
  \item perform the \(\kappa\)-framework fit using SMEFiT.
\end{enumerate}
This is not yet a full detector-level \LHmuC{} Higgs global fit.

\subsection{Vector, fermion, and loop modifier plots}

\begin{figure}[H]
\centering
\includefig[width=0.72\linewidth]{LHmuC_vector.pdf}
\caption{Vector-modifier uncertainties in the \(\kappa\)-3 framework.  The clearest
improvements appear in \(\kappa_W\) and \(\kappa_Z\), especially for higher-energy
\LHmuC{} stages.}
\label{fig:kappa_vector}
\end{figure}

\begin{figure}[H]
\centering
\begin{subfigure}{0.48\linewidth}
\centering
\includefig[width=\linewidth]{LHmuC_fermion.pdf}
\caption{Fermion modifiers.}
\end{subfigure}
\hfill
\begin{subfigure}{0.48\linewidth}
\centering
\includefig[width=\linewidth]{LHmuC_loop.pdf}
\caption{Loop-induced modifiers.}
\end{subfigure}
\caption{Fermion and loop-induced modifier projections.  In the displayed setup, the most
visible \LHmuC{} improvement is in the vector sector.}
\label{fig:kappa_fermion_loop}
\end{figure}

The vector panel shows that \(\kappa_W\) improves from about \(0.33\%\) for the LHeC-only
case to about \(0.29\%\), \(0.19\%\), and \(0.16\%\) for LHeC+\LHmuC{}-500,
LHeC+\LHmuC{}-1000, and LHeC+\LHmuC{}-2000 respectively.  The \(\kappa_Z\) uncertainty
also improves.  The physics reason is simple: \LHmuC{} Higgs production is mostly
weak-boson fusion.

\begin{keypoint}{Main message of the global-fit slide}
\vspace{0.25cm}
Because Higgs production at a lepton--hadron collider is driven by weak-boson fusion,
\LHmuC{} most directly improves the vector Higgs couplings, especially \(\kappa_W\) and
\(\kappa_Z\).  The result is promising, but should be called an exploratory projection.
\end{keypoint}



\section{Why \LHmuC{} is a distinctive BSM machine}

\subsection{The basic mechanism}

The defining BSM advantage of \LHmuC{} is access to direct partonic initial states
\begin{equation}
  \mu q,\qquad \mu g,\qquad \gamma\gamma,\qquad \gamma q,
\end{equation}
at multi-TeV energy.  The most important resonant channels are
\begin{equation}
  \mu q\to X,\qquad \mu g\to X.
\end{equation}

At a \(pp\) collider, many new colored particles are searched for through pair production,
\[
  pp\to X\bar X,
\]
or through deviations in high-mass tails.  At \LHmuC{}, if a new particle couples directly
to a muon and a quark or gluon, it can be produced singly.  This can be a major advantage
because single production has a lower threshold than pair production.

\subsection{Natural BSM examples}

\begin{itemize}[leftmargin=*]
  \item RPV squarks: \(\mu q\to\tilde q\);
  \item leptoquarks: \(\mu q\to{\rm LQ}\);
  \item color-octet muons / leptogluons: \(\mu g\to\mu_8\);
  \item excited muons;
  \item contact interactions;
  \item photon-induced anomalous dipole and quartic-gauge-coupling probes.
\end{itemize}

\begin{keypoint}{Best sentence for the BSM motivation slide}
\vspace{0.25cm}
\LHmuC{} is not just a higher-energy DIS machine; it is a direct \(\mu q/\mu g\)
resonance machine for new physics coupled to second-generation leptons.
\end{keypoint}

\section{RPV squarks}

\subsection{R-parity violation}

If R-parity is violated, the MSSM superpotential may contain
\begin{equation}
  W_{\rm RPV}
  =
  \epsilon_iL_i\cdot H_u
  +
  \frac12\lambda_{ijk}L_i\cdot L_j\bar E_k
  +
  \lambda'_{ijk}L_i\cdot Q_j\bar D_k
  +
  \frac12\lambda''_{ijk}\bar U_i\bar D_j\bar D_k.
  \label{eq:wrpv}
\end{equation}
The \(LQD\)-type term,
\[
  \lambda'_{ijk}L_iQ_j\bar D_k,
\]
is directly relevant for lepton--hadron colliders because it couples a lepton, a quark,
and a squark.  For \(i=2\), the lepton is a muon-type lepton.

\subsection{Benchmark channels}

The benchmark processes can be written schematically as
\begin{align}
  p\mu^+ &\to d\,\mu^+ && \text{neutral-current-like channel},\\
  p\mu^+ &\to \bar d\,\bar\nu_\mu && \text{charged-current-like channel}.
\end{align}
The neutral-current channel can be stronger because it can involve valence-quark PDFs and
contains a high-\(p_T\) muon in the final state.

\subsection{How the exclusion curves are obtained}

The workflow is:
\begin{enumerate}[leftmargin=*]
  \item turn on one RPV coupling at a time;
  \item scan the squark mass and the coupling;
  \item generate signal and background with MG5\_aMC@NLO using an RPV-MSSM UFO model;
  \item shower/hadronize with Pythia6;
  \item cluster jets with FastJet;
  \item simulate detector effects with Delphes;
  \item apply event selection, including at least one jet and exactly one muon in the NC case;
  \item apply a hard \(p_T\)-sum threshold,
  \[
    p_T^{j_1}+p_T^{j_2}>p_T^{\rm cut};
  \]
  \item compute
  \[
    S_{\rm stat}=\frac{N_s}{\sqrt{N_b}};
  \]
  \item define the \(95\%\) CL exclusion by \(S_{\rm stat}=2\).
\end{enumerate}

\begin{figure}[H]
\centering
\includefig[width=0.80\linewidth]{exclusion_limits.jpg}
\caption{Representative \(95\%\) CL exclusion limits for an RPV neutral-current benchmark.
The horizontal axis is the squark mass, while the vertical axis is the RPV coupling.
Lower curves mean stronger sensitivity.  Different colored curves correspond to different
jet-\(p_T\)-sum thresholds.  The red curves show LHC and HL-LHC bounds.}
\label{fig:rpv_exclusion}
\end{figure}

The green curve uses no hard \(p_T\)-sum threshold.  It is strongest at low mass but weakens
at high mass.  The blue and purple curves impose harder cuts.  They are weaker at low mass
but better for heavier squarks, whose decay products are naturally harder.

\begin{keypoint}{RPV takeaway}
\vspace{0.25cm}
In favorable neutral-current benchmark channels, \LHmuC{} can probe
\[
  \lambda'_{2j1}\sim0.01{-}0.02
\]
for squark masses around \(1{-}3\tev\).  This can approach an order-of-magnitude improvement
over LHC bounds in the on-shell region, but it is a benchmark-specific statement.
\end{keypoint}

\section{Color-octet muons / leptogluons}

\subsection{What is \(\mu_8\)?}

A color-octet muon, denoted \(\mu_8\), is a hypothetical strongly interacting partner of
the muon.  It appears in composite-lepton scenarios, where Standard Model leptons may
have more fundamental constituents.  The SM muon is a color singlet, while \(\mu_8\)
transforms as an octet under \(SU(3)_c\).

\subsection{Effective interaction}

A common effective interaction is
\begin{equation}
  \mathcal{L}
  =
  \frac{1}{2\Lambda}
  \left\{
  \bar l_8^\alpha
  g_sG_{\mu\nu}^\alpha
  \sigma^{\mu\nu}
  \left(\eta_Ll_L+\eta_Rl_R\right)
  +{\rm h.c.}
  \right\}.
  \label{eq:leptogluon}
\end{equation}
Here \(\Lambda\) is the compositeness scale, \(g_s\) is the strong coupling,
\(G_{\mu\nu}^\alpha\) is the gluon field strength, \(l_8^\alpha\) is the color-octet lepton,
and \(\eta_L,\eta_R\) are chirality coefficients.  The benchmark usually assumes
\[
  \Lambda=m_{\mu_8}.
\]

\subsection{Why \(\mu p\) is special}

At \(pp\) colliders, \(\mu_8\) searches mainly rely on pair production:
\[
  pp\to\mu_8\bar\mu_8.
\]
At a \(\mu p\) collider, the incoming muon can scatter directly on a gluon:
\[
  \mu g\to\mu_8\to\mu g.
\]
This single-resonance channel benefits from the gluon PDF, has a lower mass threshold than
pair production, and gives a clean \(\mu+j\) final state.

\subsection{Analysis strategy}

The benchmark analysis uses:
\begin{itemize}[leftmargin=*]
  \item signal: \(\mu p\to\mu_8\to\mu j\);
  \item background: SM \(\mu p\to\mu j\);
  \item exactly one muon and one jet;
  \item \(p_T^\mu>400\gev\), \(p_T^j>400\gev\);
  \item \(|\eta_j|<3.5\), \(|\eta_\mu|<3\);
  \item extraction of signal and background from the \(m_{\mu j}\) spectrum;
  \item significance
  \[
    Z=\frac{S}{\sqrt{B}}.
  \]
\end{itemize}

\begin{figure}[H]
\centering
\includefig[width=0.82\linewidth]{mu8.jpg}
\caption{Representative statistical significance for color-octet muon searches.
The horizontal axis is \(m_{\mu_8}\); the vertical axis is \(S/\sqrt{B}\).
The horizontal lines mark \(2\sigma\), \(3\sigma\), and \(5\sigma\).  The \LHmuC{}
curves extend to higher mass than the HL-LHC benchmark because \(\mu g\to\mu_8\)
is a direct single-resonance channel.}
\label{fig:mu8_significance}
\end{figure}

The representative \(5\sigma\) reaches are approximately
\[
  m_{\mu_8}^{5\sigma}\sim
  2.3\tev\quad(\HL),
  \qquad
  3.0\tev\quad(\LHmuC\text{-}500),
  \qquad
  4.2\tev\quad(\LHmuC\text{-}1000).
\]

\begin{keypoint}{Best slide sentence}
\vspace{0.25cm}
The key advantage is the resonant \(\mu g\) channel,
\[
  \mu g\to\mu_8\to\mu g,
\]
which is absent in \(pp\) pair-production searches.
\end{keypoint}

\section{How to explain each plot}

\subsection{\(h\to b\bar b\) table}

The table summarizes signal and background cross sections and the final projected rate and
coupling sensitivities.  The key point is not only that the signal rate is large, but that the
backgrounds can be suppressed enough by the two-\(b\)-jet requirement and the Higgs mass
window to make a percent-level rate measurement plausible.

Suggested caption:
\begin{quote}
Cut-based charged-current \(h\to b\bar b\) benchmark.  The quoted \(g_{hbb}\) reach follows
from the measured \(j\bar\nu_\mu b\bar b\) rate after propagating external information on
\(g_{hWW}^2/\Gamma_h\).
\end{quote}

\subsection{Vector \(\kappa\)-modifier plot}

The vector plot shows projected uncertainties on \(\kappa_W\), \(\kappa_Z\), \({\rm Br}_{\rm inv}\),
and \({\rm Br}_{\rm und}\).  The strongest trend is the improvement in \(\kappa_W\) and
\(\kappa_Z\).  The reason is that \LHmuC{} Higgs production is dominated by weak-boson fusion.

Suggested caption:
\begin{quote}
Exploratory \(\kappa\)-framework projection using SMEFiT.  \LHmuC{} improves the vector
modifiers most strongly, under LHeC-like systematic assumptions and \(\kappa_V\le1\).
\end{quote}

\subsection{RPV plot}

The RPV plot shows \(95\%\) CL exclusion curves in the plane of squark mass and RPV coupling.
Lower curves are better.  Harder \(p_T\)-sum cuts are useful at high mass because heavy
resonances produce harder jets.

Suggested caption:
\begin{quote}
RPV squark coupling reach from \(\mu q\)-initiated resonance production.  Different curves
show different jet-\(p_T\)-sum thresholds; lower curves correspond to stronger coupling reach.
\end{quote}

\subsection{\(\mu_8\) plot}

The \(\mu_8\) plot shows \(S/\sqrt{B}\) as a function of \(m_{\mu_8}\).  The \LHmuC{} curves
remain above the \(5\sigma\) line to larger masses than the HL-LHC benchmark.  This is the
visual demonstration of the single-resonance \(\mu g\) advantage.

Suggested caption:
\begin{quote}
Leptogluon reach from resonant \(\mu g\to\mu_8\to\mu g\).  The \LHmuC{} curves extend to
higher masses than the HL-LHC pair-production benchmark.
\end{quote}

\section{Improved slide text}

\subsection{Slide: concrete Higgs benchmark}

\textbf{Improved title:} Concrete Higgs benchmark: \(h\to b\bar b\) in charged-current VBF

\textbf{Improved bullets:}
\begin{itemize}[leftmargin=*]
  \item Benchmark channel:
  \[
    p\mu^+\to j\bar\nu_\mu h,\qquad h\to b\bar b.
  \]
  \item Simulation chain: MG5\_aMC@NLO + Pythia6 + FastJet + Delphes.
  \item LHeC-like asymmetric detector setup with \(75\%\) \(b\)-tagging.
  \item Selection: exactly two \(b\)-jets and \(|m_{bb}-m_h|<25\gev\).
  \item Rate interpretation:
  \[
    \sigma(j\bar\nu_\mu b\bar b)\propto
    \frac{g_{hWW}^2g_{hbb}^2}{\Gamma_h}.
  \]
\end{itemize}

\textbf{Key finding:}
Already with LHeC-like detector assumptions, \LHmuC{} reaches a sub-percent \(g_{hbb}\)
benchmark sensitivity.  Final robustness requires validated luminosity, detector/MDI
background, and systematic assumptions.

\subsection{Slide: global Higgs fit}

\textbf{Improved title:} Exploratory global Higgs fit: vector modifiers benefit most

\textbf{Improved bullets:}
\begin{itemize}[leftmargin=*]
  \item Beyond the dedicated \(h\to b\bar b\) benchmark, we perform an exploratory
        \(\kappa\)-framework fit.
  \item LHeC systematic scale factors are kept fixed.
  \item Statistical uncertainties are rescaled using \LHmuC{} benchmark Higgs rates.
  \item The fit is performed with SMEFiT.
  \item The clearest projected improvement is in the vector sector:
  \[
    \delta\kappa_W\simeq0.29\%~(\LHmuC\text{-}500),
    \qquad
    0.19\%~(\LHmuC\text{-}1000).
  \]
\end{itemize}

\textbf{Caveat:}
The fit assumes \(\kappa_V\le1\) to control the Higgs-width flat direction; this is a standard
benchmark assumption, not a fully model-independent statement.

\subsection{Slide: why \LHmuC{} is a BSM machine}

\textbf{Improved title:} Why \LHmuC{} is a distinctive BSM machine

\textbf{Improved bullets:}
\begin{itemize}[leftmargin=*]
  \item \LHmuC{} combines a point-like muon beam with quark and gluon partons at multi-TeV energy.
  \item Its channel-specific advantage is direct access to
  \[
    \mu q\to X,\qquad \mu g\to X.
  \]
  \item This enables resonant single production of new states coupled to second-generation leptons.
  \item Representative portals: RPV squarks, leptoquarks, color-octet muons, excited leptons,
        and contact interactions.
  \item Compared with \(pp\), \LHmuC{} can test selected models through single-resonance production
        rather than only pair production or inclusive tails.
\end{itemize}

\subsection{Slide: representative BSM benchmarks}

\textbf{Improved title:} Benchmark BSM portals: \(\mu q\) squarks and \(\mu g\) leptogluons

\textbf{Improved bullets:}
\begin{itemize}[leftmargin=*]
  \item \textbf{RPV SUSY:} \(LQD\)-type couplings give resonant squark production in
        \(\mu q\) scattering.  In favorable NC channels,
  \[
    \lambda'_{2j1}\sim0.01{-}0.02
  \]
  can be probed for \(m_{\tilde q}\sim1{-}3\tev\).
  \item \textbf{Color-octet muon / leptogluon:}
  \[
    \mu g\to\mu_8\to\mu g.
  \]
  This direct \(\mu g\) resonance is absent in \(pp\) pair-production searches.
  \item Representative \(5\sigma\) reach:
  \[
    \HL:\ 2.3\tev,\qquad
    \LHmuC\text{-}500:\ 3.0\tev,\qquad
    \LHmuC\text{-}1000:\ 4.2\tev.
  \]
\end{itemize}

\textbf{Caveat:}
Representative benchmark studies; final reach requires detector/MDI backgrounds, systematic
uncertainties, and channel-specific optimization.

\section{Speaker notes}

\subsection{\(h\to b\bar b\)}

This is the most concrete Higgs benchmark in the talk.  The process is charged-current
VBF, \(p\mu^+\to j\bar\nu_\mu h\), followed by the dominant Higgs decay \(h\to b\bar b\).
The signal is selected with two \(b\)-jets and a Higgs mass window.  The measured rate scales as
\(g_{hWW}^2g_{hbb}^2/\Gamma_h\), so after using external information on
\(g_{hWW}^2/\Gamma_h\), we obtain a projected \(g_{hbb}\) sensitivity.  The key point is that
the large \LHmuC{} rate turns this into a sub-percent benchmark, although the final number
depends on luminosity and systematic assumptions.

\subsection{Global fit}

The previous slide was a single-channel study.  This slide asks what happens in a broader Higgs
coupling interpretation.  We use a \(\kappa\)-framework fit, implemented with SMEFiT.  This is
exploratory: we keep the LHeC systematic scale factors and rescale the statistical part using the
\LHmuC{} rates.  The clearest improvement is in the vector modifiers, especially \(\kappa_W\) and
\(\kappa_Z\), which is natural because Higgs production at a lepton--hadron collider is primarily
weak-boson fusion.

\subsection{BSM motivation}

The BSM point is simple but important.  In \(pp\), many colored new particles are searched for
through pair production, which costs energy because two heavy particles must be produced.  At
\LHmuC{}, if the new state couples to a muon and a quark or gluon, it can be produced singly as
a resonance.  This is why RPV squarks and color-octet muons are natural benchmark examples.

\subsection{BSM benchmarks}

These two examples show the same BSM principle in two different initial states.  RPV squarks use
\(\mu q\) resonance production through \(LQD\)-type couplings, and favorable NC channels can
probe \(\lambda'\) around \(10^{-2}\).  Color-octet muons use the \(\mu g\) initial state.  The
single-resonance channel \(\mu g\to\mu_8\to\mu g\) gives a stronger reach than HL-LHC pair
production in this benchmark.  The exact numbers should be understood as representative, pending
full detector and MDI studies.

\section{Likely audience questions and compact answers}

\paragraph{Q: How robust is the \(hbb\) sensitivity once systematics are included?}
The quoted rate precision is primarily statistical.  The coupling extraction includes external
information on \(g_{hWW}^2/\Gamma_h\), but a final study must include \(b\)-tagging, jet-energy
scale, background normalization, theory/PDF uncertainties, and beam-induced backgrounds.

\paragraph{Q: Why do you need external information on \(g_{hWW}^2/\Gamma_h\)?}
Because the measured rate is proportional to \(g_{hWW}^2g_{hbb}^2/\Gamma_h\).  A single rate
cannot isolate \(g_{hbb}\) without information on the production coupling and total width
combination.

\paragraph{Q: Is the global fit independent, or mostly a rescaling?}
It should be presented as an exploratory projection.  Statistical uncertainties are rescaled using
\LHmuC{} cross sections, while LHeC systematic factors are retained.

\paragraph{Q: Why impose \(\kappa_V\le1\)?}
It controls the Higgs-width flat direction in fits without a direct model-independent Higgs width
measurement.  It is useful for comparison, but not fully model independent.

\paragraph{Q: Why are RPV squarks like leptoquarks here?}
The \(LQD\) operator couples a lepton, a quark, and a squark.  Phenomenologically, the squark
can appear as a leptoquark-like resonance in \(\mu q\) scattering.

\paragraph{Q: Are interference effects included in the RPV benchmark?}
The benchmark treatment neglects subdominant interference between signal and background.  A
final study should include it.

\paragraph{Q: Why is \(\mu p\) better than \(pp\) for \(\mu_8\)?}
Because \(\mu p\) provides the direct single-resonance process \(\mu g\to\mu_8\), whereas \(pp\)
searches mainly rely on pair production.

\paragraph{Q: How model-dependent is the \(\mu_8\) reach?}
It depends strongly on the effective operator, the chirality assumptions, and especially \(\Lambda\).
The benchmark assumes \(\Lambda=m_{\mu_8}\).

\paragraph{Q: What about muon-decay beam-induced backgrounds?}
This is a central detector/MDI issue.  The benchmark studies are promising, but realistic shielding,
timing, and detector-response simulations are needed before final experimental claims.

\section{Final checklist before the presentation}

\subsection*{Must fix}
\begin{enumerate}[leftmargin=*]
  \item Resolve the \(hbb\) luminosity consistency issue for the \(1\tev\) benchmark.
  \item Label \(\sigma_{\rm sig}\) clearly: inclusive \(h\) production or final \(b\bar b\) state.
  \item Replace ``explorative'' with ``exploratory''.
  \item Add one caveat line to the global-fit slide: LHeC-like systematic factors are retained.
  \item Fix spacing: ``\LHmuC{} benchmark'', not ``\LHmuC{}benchmark''.
\end{enumerate}

\subsection*{Should improve if time}
\begin{enumerate}[leftmargin=*]
  \item Add a short explanation of \(\kappa_V\le1\).
  \item Mention \(\Lambda=m_{\mu_8}\) in the \(\mu_8\) caption or speaker note.
  \item Make plot captions carry the physics message.
  \item Replace ``Color-octet muon \(\mu_8\) reach'' by
  ``Leptogluon reach from resonant \(\mu g\to\mu_8\to\mu g\)''.
\end{enumerate}

\section{Further reading map}

\begin{longtable}{@{}p{0.23\linewidth}p{0.35\linewidth}p{0.34\linewidth}@{}}
\caption{Suggested references and what to learn from each.}\\
\toprule
Topic & Reference & Why read it \\
\midrule
\endfirsthead
\toprule
Topic & Reference & Why read it \\
\midrule
\endhead
\LHmuC{} concept and HELP4CERN physics case &
HELP4CERN white paper draft &
Overall collider motivation, benchmark cross sections, Higgs, \(\gamma\gamma\), top, and BSM overview. \\
Muon--proton Higgs and RPV benchmarks &
Cheung and Wang, Phys. Rev. D 103, 116009 (2021), arXiv:2101.10476 &
Original \(\mu p\) benchmark for \(hbb\) and RPV squark reach. \\
Future Higgs \(\kappa\)-framework &
de Blas et al., JHEP 01 (2020) 139, arXiv:1905.03764 &
Standard future-collider Higgs comparison and \(\kappa\)-framework context. \\
Future collider SMEFT fits &
de Blas et al., arXiv:2206.08326 &
SMEFT global-fit context and future-collider projections. \\
SMEFiT and precision reach &
Armadillo et al., arXiv:2604.16596 &
Recent multi-framework precision interpretation and updated SMEFiT release. \\
\(\mu\)LHC and color-octet muons &
Akturk et al., arXiv:2506.15445 &
\(\mu^+p\) collider design and color-octet muon benchmark. \\
Historical RPV at \(\mu p\) &
Carena, Choudhury, Quigg, Raychaudhuri, Phys. Rev. D 62, 095010 &
Classic study of RPV resonances at muon--proton colliders. \\
Leptogluon phenomenology &
Goncalves-Netto et al.; Acar--Kaya--Oner &
Color-octet lepton / leptogluon theory and collider phenomenology. \\
\bottomrule
\end{longtable}

\begin{thebibliography}{99}

\bibitem{HELP4CERN}
D.~Acosta et al.,
\emph{White Paper on High Energy, High Luminosity Lepton--Hadron Colliders at CERN},
HELP4CERN Study Group draft.

\bibitem{KhanpourDIS}
H.~Khanpour,
\emph{Perspectives for HET physics studies at a muon-hadron collider at CERN},
DIS 2026 presentation draft.

\bibitem{CheungWang}
K.~Cheung and Z.~S.~Wang,
\emph{Physics potential of a muon-proton collider},
Phys. Rev. D \textbf{103}, 116009 (2021),
doi:10.1103/PhysRevD.103.116009,
arXiv:2101.10476.

\bibitem{deBlasHiggs}
J.~de Blas et al.,
\emph{Higgs Boson Studies at Future Particle Colliders},
JHEP \textbf{01} (2020) 139,
arXiv:1905.03764.

\bibitem{deBlasSMEFT}
J.~de Blas, Y.~Du, C.~Grojean, J.~Gu, V.~Miralles, M.~E.~Peskin, J.~Tian, M.~Vos and E.~Vryonidou,
\emph{Global SMEFT Fits at Future Colliders},
Snowmass 2021,
arXiv:2206.08326.

\bibitem{Armadillo}
T.~Armadillo, E.~Celada, J.~ter Hoeve, F.~Maltoni, L.~Mantani, J.~Rojo, A.~N.~Rossia,
S.~Tentori, M.~O.~A.~Thomas and E.~Vryonidou,
\emph{New Physics Reach through Precision at Future Colliders: a Multi-Pronged Approach},
arXiv:2604.16596.

\bibitem{Akturk}
D.~Akturk et al.,
\emph{\(\mu\)LHC: Antimuon Ring and HL-LHC based \(\mu^+p\) Collider},
arXiv:2506.15445.

\bibitem{Carena}
M.~Carena, D.~Choudhury, C.~Quigg and S.~Raychaudhuri,
\emph{Study of R-parity violation at a muon proton collider},
Phys. Rev. D \textbf{62}, 095010 (2000),
arXiv:hep-ph/0006144.

\bibitem{Leptogluon}
D.~Goncalves-Netto, D.~Lopez-Val, K.~Mawatari, I.~Wigmore and T.~Plehn,
\emph{Looking for leptogluons},
Phys. Rev. D \textbf{87} (2013).

\bibitem{Acar}
Y.~C.~Acar, U.~Kaya and B.~B.~Oner,
\emph{Resonant production of color octet muons at Future Circular Collider-based muon-proton colliders},
Chin. Phys. C \textbf{42} (2018) 083108.

\bibitem{MG5}
J.~Alwall et al.,
\emph{MadGraph 5: Going Beyond},
JHEP \textbf{06} (2011) 128.

\bibitem{FastJet}
M.~Cacciari, G.~P.~Salam and G.~Soyez,
\emph{FastJet user manual},
Eur. Phys. J. C \textbf{72} (2012) 1896.

\bibitem{Delphes}
J.~de Favereau et al.,
\emph{DELPHES 3: a modular framework for fast simulation of a generic collider experiment},
JHEP \textbf{02} (2014) 057.

\end{thebibliography}

\end{document}
